% Figures for ITP 2026 Paper
% These are TikZ figures to be included in paper.tex

% Required packages (add to preamble):
% \usepackage{tikz}
% \usetikzlibrary{shapes,arrows,positioning,fit,backgrounds,calc}

%%%%%%%%%%%%%%%%%%%%%%%%%%%%%%%%%%%%%%%%%%%%%%%%%%%%%%%%%%%%%%%%%%%%%%%%%%%%%%%
% FIGURE 1: Game State Structure
%%%%%%%%%%%%%%%%%%%%%%%%%%%%%%%%%%%%%%%%%%%%%%%%%%%%%%%%%%%%%%%%%%%%%%%%%%%%%%%

% Shows the hierarchical structure of GameState with zones

\begin{figure}[t]
\centering
\begin{tikzpicture}[
    node distance=0.3cm,
    zone/.style={draw, rounded corners, minimum width=1.8cm, minimum height=0.6cm, font=\small},
    player/.style={draw, rounded corners, thick, inner sep=0.3cm},
    label/.style={font=\footnotesize\itshape},
    >=stealth
]

% Player One
\node[zone, fill=blue!10] (deck1) {Deck};
\node[zone, fill=blue!10, right=of deck1] (hand1) {Hand};
\node[zone, fill=blue!10, right=of hand1] (prizes1) {Prizes};
\node[zone, fill=blue!20, below=0.4cm of hand1] (active1) {Active};
\node[zone, fill=blue!15, below=0.2cm of active1, minimum width=4cm] (bench1) {Bench (0--5)};
\node[zone, fill=blue!10, below=0.2cm of bench1] (discard1) {Discard};

\node[player, fit=(deck1)(hand1)(prizes1)(active1)(bench1)(discard1), label=above:{\textbf{Player One}}] (p1) {};

% Player Two (mirror)
\node[zone, fill=red!10, right=2.5cm of prizes1] (deck2) {Deck};
\node[zone, fill=red!10, right=of deck2] (hand2) {Hand};
\node[zone, fill=red!10, right=of hand2] (prizes2) {Prizes};
\node[zone, fill=red!20, below=0.4cm of hand2] (active2) {Active};
\node[zone, fill=red!15, below=0.2cm of active2, minimum width=4cm] (bench2) {Bench (0--5)};
\node[zone, fill=red!10, below=0.2cm of bench2] (discard2) {Discard};

\node[player, fit=(deck2)(hand2)(prizes2)(active2)(bench2)(discard2), label=above:{\textbf{Player Two}}] (p2) {};

% Active player indicator
\node[below=0.8cm of $(bench1)!0.5!(bench2)$, font=\small] (api) {activePlayer : PlayerId};

% Battle arrow
\draw[<->, thick, red!60!black] (active1.east) -- node[above, font=\scriptsize] {Attack} (active2.west);

\end{tikzpicture}
\caption{Game state structure. Each player maintains six zones; the active Pok\'emon battle each other.}
\label{fig:gamestate}
\end{figure}


%%%%%%%%%%%%%%%%%%%%%%%%%%%%%%%%%%%%%%%%%%%%%%%%%%%%%%%%%%%%%%%%%%%%%%%%%%%%%%%
% FIGURE 2: Damage Calculation Pipeline  
%%%%%%%%%%%%%%%%%%%%%%%%%%%%%%%%%%%%%%%%%%%%%%%%%%%%%%%%%%%%%%%%%%%%%%%%%%%%%%%

\begin{figure}[t]
\centering
\begin{tikzpicture}[
    node distance=0.8cm,
    block/.style={draw, rounded corners, minimum width=2.2cm, minimum height=0.8cm, align=center, font=\small},
    arrow/.style={->, thick, >=stealth},
    op/.style={font=\footnotesize, fill=white, inner sep=2pt}
]

% Pipeline stages
\node[block, fill=yellow!20] (base) {Base\\Damage};
\node[block, fill=orange!20, right=of base] (effects) {+ Effects\\Bonus};
\node[block, fill=red!20, right=of effects] (weak) {$\times$ Weakness\\(2 if match)};
\node[block, fill=green!20, right=of weak] (resist) {$-$ Resistance\\(30 if match)};
\node[block, fill=blue!20, right=of resist] (final) {Final\\Damage};

% Arrows
\draw[arrow] (base) -- (effects);
\draw[arrow] (effects) -- (weak);
\draw[arrow] (weak) -- (resist);
\draw[arrow] (resist) -- (final);

% Example values below
\node[below=0.5cm of base, font=\scriptsize] {20};
\node[below=0.5cm of effects, font=\scriptsize] {20 + 0};
\node[below=0.5cm of weak, font=\scriptsize] {20 $\times$ 2};
\node[below=0.5cm of resist, font=\scriptsize] {40 $-$ 0};
\node[below=0.5cm of final, font=\scriptsize\bfseries] {= 40};

% Label
\node[below=1.2cm of $(base)!0.5!(final)$, font=\footnotesize\itshape] {Example: Squirtle (Water) vs Charmander (weak to Water)};

\end{tikzpicture}
\caption{Damage calculation pipeline with weakness/resistance modifiers.}
\label{fig:damage}
\end{figure}


%%%%%%%%%%%%%%%%%%%%%%%%%%%%%%%%%%%%%%%%%%%%%%%%%%%%%%%%%%%%%%%%%%%%%%%%%%%%%%%
% FIGURE 3: Theorem Architecture
%%%%%%%%%%%%%%%%%%%%%%%%%%%%%%%%%%%%%%%%%%%%%%%%%%%%%%%%%%%%%%%%%%%%%%%%%%%%%%%

\begin{figure}[t]
\centering
\begin{tikzpicture}[
    node distance=0.6cm and 0.8cm,
    thm/.style={draw, rounded corners, minimum width=2.8cm, minimum height=0.6cm, align=center, font=\scriptsize},
    cat/.style={font=\footnotesize\bfseries, text=gray},
    arrow/.style={->, >=stealth, thick, gray},
    >=stealth
]

% Semantics layer
\node[cat] (sem) {Semantics};
\node[thm, fill=blue!15, below=0.3cm of sem] (legal) {Legal\_iff\_preconditions};
\node[thm, fill=blue!15, right=of legal] (total) {step\_total\_for\_legal};
\node[thm, fill=blue!15, left=of legal] (error) {step\_error\_iff\_not\_legal};

% Invariants layer
\node[cat, below=1cm of legal] (inv) {Invariants};
\node[thm, fill=green!15, below=0.3cm of inv] (valid) {step\_preserves\_valid};
\node[thm, fill=green!15, right=of valid] (cards) {step\_preserves\_total\_cards};
\node[thm, fill=green!15, left=of valid] (won) {step\_preserves\_hasWon};

% Reachability layer  
\node[cat, below=1cm of valid] (reach) {Reachability};
\node[thm, fill=yellow!20, below=0.3cm of reach] (rvalid) {reachable\_valid\_initial};
\node[thm, fill=yellow!20, right=of rvalid] (rcons) {reachable\_card\_conservation};
\node[thm, fill=yellow!20, left=of rvalid] (nt1w) {no\_turn\_one\_win};

% Solver layer
\node[cat, below=1cm of rvalid] (solv) {Solver};
\node[thm, fill=red!15, below=0.3cm of solv] (sound) {bestAttack\_sound};
\node[thm, fill=red!15, right=of sound] (opt) {bestAttack\_optimal};
\node[thm, fill=red!15, left=of sound] (twoply) {solveTwoPly\_optimal};

% Dependency arrows
\draw[arrow] (legal) -- (total);
\draw[arrow] (total) -- (valid);
\draw[arrow] (valid) -- (rvalid);
\draw[arrow] (cards) -- (rcons);
\draw[arrow] (legal) -- (sound);

\end{tikzpicture}
\caption{Theorem architecture: semantics foundations support invariant preservation, enabling reachability and solver proofs.}
\label{fig:theorems}
\end{figure}


%%%%%%%%%%%%%%%%%%%%%%%%%%%%%%%%%%%%%%%%%%%%%%%%%%%%%%%%%%%%%%%%%%%%%%%%%%%%%%%
% FIGURE 4: Case Study - Solver Decision
%%%%%%%%%%%%%%%%%%%%%%%%%%%%%%%%%%%%%%%%%%%%%%%%%%%%%%%%%%%%%%%%%%%%%%%%%%%%%%%

\begin{figure}[t]
\centering
\begin{tikzpicture}[
    pokemon/.style={draw, rounded corners, minimum width=2.5cm, minimum height=1.5cm, align=center},
    stat/.style={font=\scriptsize},
    arrow/.style={->, thick, >=stealth},
    result/.style={draw, dashed, rounded corners, fill=gray!10, minimum width=2cm, align=center, font=\small}
]

% Squirtle
\node[pokemon, fill=blue!20] (squirtle) {
    \textbf{Squirtle}\\
    \stat{Type: Water}\\
    \stat{HP: 60}\\
    \stat{Energy: [W]}
};
\node[below=0.1cm of squirtle, font=\scriptsize] {Attack: Water Gun (20)};

% Charmander
\node[pokemon, fill=red!20, right=3cm of squirtle] (charmander) {
    \textbf{Charmander}\\
    \stat{Type: Fire}\\
    \stat{HP: 70, Dmg: 0}\\
    \stat{Weak: Water}
};

% Arrow with damage
\draw[arrow, blue!60!black, very thick] (squirtle.east) -- node[above, font=\small] {40 damage} node[below, font=\scriptsize] {(20 $\times$ 2 weakness)} (charmander.west);

% Solver result box
\node[result, below=1cm of $(squirtle)!0.5!(charmander)$] (result) {
    \textbf{Solver Output}\\
    \stat{attackIndex: 0}\\
    \stat{damage: 40}\\
    \stat{isLethal: false}
};

% Verified badge
\node[right=0.3cm of result, font=\scriptsize, text=green!50!black] {$\checkmark$ Proven optimal};

\end{tikzpicture}
\caption{Case study: the certified solver selects Water Gun against Charmander, correctly applying the $\times 2$ weakness multiplier.}
\label{fig:casestudy}
\end{figure}
